%%%
%%% Radical "⾹"
%%%
\section*{Radical 186: ``⾹''}\addcontentsline{toc}{section}{Radical 186: ⾹}\addcontentsline{loh}{figure}{\#\#\#\# 186: ⾹}

%%%%%%%%%% 香 %%%%%%%%%%
\subsection*{香}\addcontentsline{loh}{figure}{香}

\begin{Entry}{香}{9}{⾹}[Kangxi 186]
  \begin{Phonetics}{香}{xiang1}[][HSK 3]
    \definition*{s.}{Sobrenome: Xiang}
    \definition{adj.}{aromático; perfumado; fragrante; cheiroso | saboroso; saboroso; delicioso; apetitoso | com gosto; com bom apetite | (sono) profundo; dormir confortavelmente e tranquilamente | popular; valorizado; apreciado}
    \definition[根,炷]{s.}{especiaria; perfume; fragrância; aromatizante; substância com aroma intenso | incenso; bastão de incenso; tiras finas feitas de serragem e especiarias, queimadas em rituais para honrar os antepassados ou deuses e budas, e também usadas para afastar odores desagradáveis ou mosquitos| antigamente, referia"-se a coisas relacionadas com mulheres}
  \antonymref{臭}{chou4}
  \end{Phonetics}
\end{Entry}

\begin{Entry}{香气}{9,4}{⾹,⽓}
  \begin{Phonetics}{香气}{xiang1qi4}
    \definition{s.}{cheiro doce; fragrância; aroma | fragrância; os produtos cosméticos são feitos de fragrâncias, álcool e água destilada, entre outras coisas}
  \synonymref{香味}{xiang1wei4}
  \synonymref{直到}{zhi2dao4}
  \antonymref{臭气}{chou4qi4}
  \end{Phonetics}
\end{Entry}

\begin{Entry}{香水}{9,4}{⾹,⽔}
  \begin{Phonetics}{香水}{xiang1shui3}[][HSK 7-9]
    \definition[瓶,款,种,滴]{s.}{aroma; perfume; cosméticos feitos com fragrâncias, álcool, água destilada, etc.}
  \end{Phonetics}
\end{Entry}

\begin{Entry}{香皂}{9,7}{⾹,⽩}
  \begin{Phonetics}{香皂}{xiang1zao4}
    \definition[块,盒]{s.}{sabonete de toalete; sabonete perfumado; sabonete aromatizado; o sabonete feito com a adição de fragrância a matérias"-primas refinadas é usado principalmente para lavar o rosto}
  \synonymref{肥皂}{fei2zao4}
  \end{Phonetics}
\end{Entry}

\begin{Entry}{香肠}{9,7}{⾹,⾁}
  \begin{Phonetics}{香肠}{xiang1chang2}[][HSK 5]
    \definition[根]{s.}{salsicha; linguiça; alimento feito com intestino de porco, recheado com carne picada e temperos}
  \end{Phonetics}
\end{Entry}

\begin{Entry}{香味}{9,8}{⾹,⼝}
  \begin{Phonetics}{香味}{xiang1wei4}[][HSK 7-9]
    \definition[股,阵]{s.}{aroma; odor; cheiro doce; fragrância; o termo se refere ao conjunto de sentidos do olfato e do paladar que evoca uma sensação agradável e confortável; é percebido através dos sentidos do olfato e do paladar}
  \synonymref{香气}{xiang1qi4}
  \end{Phonetics}
\end{Entry}

\begin{Entry}{香油}{9,8}{⾹,⽔}
  \begin{Phonetics}{香油}{xiang1you2}[][HSK 7-9]
    \definition{s.}{óleo perfumado | óleo de gergelim}
  \end{Phonetics}
\end{Entry}

\begin{Entry}{香波}{9,8}{⾹,⽔}
  \begin{Phonetics}{香波}{xiang1bo1}
    \definition{s.}{Empréstimo Linguístico: xampu}
  \seealsoref{洗发皂}{xi3fa4 zao4}
  \end{Phonetics}
\end{Entry}

\begin{Entry}{香炉}{9,8}{⾹,⽕}
  \begin{Phonetics}{香炉}{xiang1lu2}
    \definition{s.}{um incensário (para queimar incenso); incensário | queimador de incenso | turíbulo}
  \end{Phonetics}
\end{Entry}

\begin{Entry}{香料}{9,10}{⾹,⽃}
  \begin{Phonetics}{香料}{xiang1liao4}[][HSK 7-9]
    \definition[种]{s.}{perfume; especiaria; condimento; substâncias orgânicas que emitem fragrância à temperatura ambiente}[这种香料可以让菜更美味。===Este tempero pode deixar os pratos mais saborosos.]
  \end{Phonetics}
\end{Entry}

\begin{Entry}{香烟}{9,10}{⾹,⽕}
  \begin{Phonetics}{香烟}{xiang1yan1}[][HSK 7-9]
    \definition[根,支,盒,包]{s.}{a fumaça da queima de incenso | cigarro}
  \synonymref{烟草}{yan1cao3}
  \synonymref{纸烟}{zhi3yan1}
  \end{Phonetics}
\end{Entry}

\begin{Entry}{香艳}{9,10}{⾹,⾊}
  \begin{Phonetics}{香艳}{xiang1yan4}
    \definition{adj.}{sexy; erótico; pornográfico | florido; de boudoir (isto é, de meninas ou mulheres); com linguagem rebuscada ou conteúdo relacionado aos aposentos femininos}
  \synonymref{艳丽}{yan4li4}
  \antonymref{平凡}{ping2fan2}
  \antonymref{朴素}{pu3su4}
  \end{Phonetics}
\end{Entry}

\begin{Entry}{香港}{9,12}{⾹,⽔}
  \begin{Phonetics}{香港}{xiang1gang3}
    \definition*{s.}{Hong Kong}
  \seealsoref{香港岛}{xiang1gang3dao3}
  \end{Phonetics}
\end{Entry}

\begin{Entry}{香港岛}{9,12,7}{⾹,⽔,⼭}
  \begin{Phonetics}{香港岛}{xiang1gang3dao3}
    \definition*{s.}{Ilha de Hong Kong}
  \seealsoref{香港}{xiang1gang3}
  \end{Phonetics}
\end{Entry}

\begin{Entry}{香槟酒}{9,14,10}{⾹,⽊,⾣}
  \begin{Phonetics}{香槟酒}{xiang1bin1jiu3}
    \definition[瓶,杯,口]{s.}{Empréstimo Linguístico: \emph{champagne}; espumante}
  \synonymref{葡萄酒}{pu2tao2jiu3}
  \end{Phonetics}
\end{Entry}

\begin{Entry}{香蕈}{9,15}{⾹,⾋}
  \begin{Phonetics}{香蕈}{xiang1xun4}
    \definition{s.}{cogumelo perfumado (\emph{Lentinus edodes}); fungo shiitake, um cogumelo comestível}
  \end{Phonetics}
\end{Entry}

\begin{Entry}{香蕉}{9,15}{⾹,⾋}
  \begin{Phonetics}{香蕉}{xiang1jiao1}[][HSK 3]
    \definition[枝,根,个,把,串,束,弓]{s.}{banana}
  \end{Phonetics}
\end{Entry}

%%%%% EOF %%%%%

